%% Interstitial: Paper 2 -> Paper 3
%% Philosophical and conceptual bridge: dynamics to computation.

Chapter~\ref{chap:paper-two} established that the predictive operator
family $\{K_t\}$ forms a semigroup and is dual to the Koopman operator.
The temporal structure of prediction is forced: given $P$, the semigroup
exists, is unique, and its algebraic properties are theorems, not
assumptions.

\medskip

\paragraph{The residual freedom.}
What the semigroup $\{K_t\}$ does not specify is the \emph{form} the
observer's queries take.  The abstract characterisation holds for any
observable dynamics satisfying the program's structural hypotheses ---
the queries are a system, the measure $P$ is canonical, the semigroup
emerges.  But nothing in Papers~1 or~2 selects the form of those
queries.  An observer measuring a physical system has finitely many
sensors, at specific locations and times; their discriminative capacity
is local and temporal.  Which query system formalises this?

A second freedom: even if the form of the queries is fixed, the minimal
sufficient window --- how much history the observer needs to retain full
predictive content --- appears open.  Retaining more history is
always safe; it is unclear what the \emph{minimum} is.

\medskip

\paragraph{The elimination.}
Chapter~\ref{chap:paper-three} shows both freedoms are illusory.

The natural query for an observer embedded in a temporal process with
local sensor access is the \emph{delay query}: $Q_{m,\tau}(\omega) =
(\omega(0), \omega(\tau), \ldots, \omega((m-1)\tau))$ records $m$
consecutive observations spaced $\tau$ apart.  These queries formalise
the Takens embedding strategy; their refinement structure is governed
by the program's directedness conditions, and the delay query system
is upper-directed --- so the extension theorem applies, and the program's
full machinery instantiates in this setting.

The minimal sufficient window is not chosen: the \emph{Markov order}
$m(\tau, P)$ is determined jointly by the measure $P$ and the lag $\tau$.
It is the minimum $m$ such that $Q_{m,\tau}$ is predictively sufficient.
The observer cannot choose to need less history than $m(\tau,P)$ and
cannot gain predictive content by taking more; the structure of $P$ and
the lag $\tau$ together determine what is necessary and what is redundant.

The operator $K_t$ is then approximated by $K_t^N$, computed from $N$
data points via local linear regression on the delay embedding.  As
$N \to \infty$, $K_t^N \to K_t$ in $L^2(P)$.  The observer, embedded
in time and accumulating experience, converges to the true predictive
structure.  The approximation is not a heuristic; it is the formal
expression of learning from within.

\medskip

\paragraph{The arc complete.}
The program began with a question about what an observer must be for
coherent probabilistic experience to be possible.  It ends with an
algorithm.  The observer's discriminative acts force a $\sigma$-algebra;
collective exhaustion forces $\sigma$-additivity; realizability and
directedness force a global probability $P$; $P$ forces the predictive
structure of every query; temporal coherence forces the semigroup; the
structure of the observer's situation forces the delay query and the
minimal sufficient window.

At each step, what appeared to be a construction is revealed to be a
constraint, and what appeared to be an assumption is revealed to be a
consequence.  The program's claim is not that it derives everything
from nothing.  It is that it derives more than was previously known to
be derivable, and identifies precisely what the remaining assumptions
are and what they mean.
